

\section{Economic application}
\label{sec:rh_econ}
\begin{itemize}
\item \textbf{A trade on the closing-bar forecast.} Our models forecast every regular-hours bar, and their
  15:30 forecast of the closing bar has a counterpart in the options market: a same-day-expiring S\&P~500 straddle
  traded at 15:30 ET and held to the close.
\end{itemize}

\subsection{Trading strategy}
\label{sec:rh_econ_strategy}
\begin{itemize}
\item \textbf{Sample.} The strategy trades on $\mathcal T_{\mathrm{opt}} \subset \mathcal T_{\mathrm{oos}}$, the
  $T_{\mathrm{opt}} = \cmNdays$ trading days of the evaluation span with a same-day SPXW expiration, early-close days
  excluded (Appendix~\ref{app:rh_econ_details}). Its first \rheSzWarmup{} days only initialize the sizing rules of
  Appendix~\ref{app:rh_econ_sizing}; every statistic of this section is computed on the remaining \rheSzDays{}
  post-warmup days, $\mathcal T_{\mathrm{opt}}^{+}$.
\item \textbf{Position.} With $S_t$ the underlying print on the option chain at 15:30 ET, we trade the call and the
  put at the nearest strikes with live quotes, $K_C = \min\{K \ge S_t\}$ and $K_P = \max\{K \le S_t\}$: an
  at-the-money straddle when $K_C = K_P$, otherwise a narrow strangle on the two strikes bracketing $S_t$, which we
  also call the straddle. It is held unhedged to cash settlement at the official S\&P~500 close
  $S^{\mathrm{close}}_t$, and its return on the mid-quote premium $P_t$, the sum of the two legs' mid-quotes, is
  \[
    R_t = \frac{\bigl(S^{\mathrm{close}}_t - K_C\bigr)^+ + \bigl(K_P - S^{\mathrm{close}}_t\bigr)^+}{P_t} - 1
      \;\ge\; -1 .
  \]
\item \textbf{Signal.} With $\mathrm{IV}^C_t$ and $\mathrm{IV}^P_t$ the legs' implied volatilities at 15:30 ET, quoted
  per hour, $\mathrm{IV}_t = (\mathrm{IV}^C_t + \mathrm{IV}^P_t)/(2\sqrt2)$ is the implied volatility over the closing
  half hour and $\mathrm{IV}_t^2$ its implied variance (Appendix~\ref{app:rh_econ_details} treats the vendor's
  failed solves). The sign strategy holds
  \[
    q_t = \sign\bigl(\hat\sigma^2_t - \mathrm{IV}_t^2\bigr), \qquad r_t = q_t R_t ,
  \]
  long when the forecast variance $\hat\sigma^2_t$ of the closing bar (Section~\ref{sec:rh_methods_setup}) exceeds
  the implied variance and short otherwise (ties go short); $\mathrm{IV}_t^2 - \hat\sigma^2_t$ is the forecast of the
  half hour's variance risk premium \citep{CarrWu2009,BollerslevTauchenZhou2009}.
\item \textbf{Benchmarks.} The unconditional short straddle, $q_t \equiv -1$, and the sign strategy of the benchmark
  forecast, the OLS on the HAR benchmark set estimated on the \rhePoolNBars{} regular-hours bars
  (Table~\ref{tab:rh_oos}). On the HAR benchmark set the ridge, lasso and elastic-net forecasts coincide with
  the OLS (Table~\ref{tab:rh_econ_main}, note b), so the six HAR-set strategies hold \rhemHarDistinctWord{} distinct
  forecasts.
\item \textbf{Execution.} At quoted prices a purchase pays the sum of the legs' asks and a sale receives the sum of
  their bids; settlement is in cash, so the spread is paid at entry only (Appendix~\ref{app:rh_econ_details}). The
  median quoted half-spread is \rhemHalfSpreadMed\% of the premium on the \rheSzDays{} post-warmup days
  (\rhemHalfSpreadMedAll\% on all \cmNdays{} days of $\mathcal T_{\mathrm{opt}}$). No commissions or fees are
  deducted.
\end{itemize}

\subsection{Performance measures}
\label{sec:rh_econ_measures}
\begin{itemize}
\item \textbf{Sharpe ratio.} With $T = \rheSzDays$ the number of post-warmup days in $\mathcal T_{\mathrm{opt}}^{+}$
  and $\mathcal Y$ their \rheNYears{} calendar years,
  \[
    \SR = \sqrt{q}\;\frac{\bar r}{\mathrm{sd}(r)},
    \qquad
    q = \frac{T}{\sum_{y \in \mathcal Y} c_y},
  \]
  where $\bar r$ and $\mathrm{sd}(r)$ are the mean and standard deviation (divisor $T - 1$) of the $T$ daily returns,
  each day with weight $1/T$, and $c_y$ is the fraction of calendar year $y$ that the sample covers, so that
  $q = \rheAnnDays$ is the average number of option days per calendar year (SPXW expirations were listed on Mondays,
  Wednesdays and Fridays until spring 2022 and on every weekday from \rheSplitDate). The ratio of a single year,
  $\SR_y = \sqrt{A_y}\,\bar r_y/\mathrm{sd}_y(r)$, takes the mean and standard deviation of the $N_y$ returns of year
  $y$ and annualizes by that year's own number of option days per calendar year, $A_y = N_y/c_y$. The risk-free
  return over the 30-minute holding period is negligible and omitted. Appendix~\ref{app:rh_econ_sharpe} gives the
  alternative definitions and when this one was adopted.
\item \textbf{Tests.} Every statement about a difference of Sharpe ratios, $\Delta\SR = \SR(r^A) - \SR(r^B)$, uses
  the studentized circular block bootstrap test of \citet{LedoitWolf2008} (Appendix~\ref{sec:rh_methods_tests}).
  A statement about
  a mean daily return, or about the mean of the return differential $\Delta_t = r^A_t - r^B_t$, uses the
  $t$-statistic with the Newey--West long-run variance at bandwidth $S = \rheHACS$ and fixed-$b$ critical values
  \citep{KieferVogelsang2005,LazarusLewisStockWatson2018}. The two tests answer different questions: a strategy can
  have a significantly higher Sharpe ratio than another without a significantly higher mean return.
\end{itemize}

\subsection{Strategy performance}
\label{sec:rh_econ_results}
\begin{itemize}
\item \textbf{Table and figure.} Table~\ref{tab:rh_econ_main} reports the sign strategies of the \rhrPoolN{}
  forecasts and the short straddle, and Figure~\ref{fig:rh_econ_byyear} plots their Sharpe ratios by calendar year,
  $\SR_y$.
\item \textbf{Passive selling earns nothing detectable.} The unconditional short straddle has $\SR = \rheShortSh$ at
  mid-quotes and \rheShortShX{} at quoted bid/ask prices (\rhemShDSixShort{} and \rhemShDOneShort{} at mid-quotes
  under the two alternative definitions of Appendix~\ref{app:rh_econ_sharpe}: the equal-weight mean of the $\SR_y$,
  which gives the partial first and last years the weight of a full one, and the $N_y$-weighted per-year ratio). Its
  mean daily return is not significantly different from zero ($t = \rheShortT$, fixed-$b$ $p = \rheShortP$), and it
  loses money in \rheShortLossYears{} of the \rheNYears{} calendar years.
\item \textbf{The benchmark forecast.} The sign strategy of the benchmark forecast reaches $\SR = \rhemShRidgeHar$ at
  mid-quotes and \rhemShXRidgeHar{} at quoted bid/ask prices. Its Sharpe ratio exceeds the short straddle's by
  \rhemDShRidgeHarVsShort{} (Ledoit--Wolf $p = \rhemPDShRidgeHarVsShort$), and none of the \rhemHarDistinctWord{}
  distinct HAR-set forecasts has a Sharpe ratio significantly above the short straddle's.
\item \textbf{The full predictor set beats passive selling.} The six forecasts on the full predictor set reach
  $\SR = \rhrPoolShFullMin$ to \rhrPoolShFullMax{} at mid-quotes and \rhrPoolShXFullMin{} to \rhrPoolShXFullMax{} at
  quoted bid/ask prices. Their Sharpe ratio gain over the short straddle is significant at the 5\% level for
  \rhemFullVsShortSigWord{} of the six at mid-quotes and \rhemFullVsShortSigXWord{} at quoted bid/ask prices
  (Ledoit--Wolf test). The mean daily return differential over the short straddle, a different test, is significant
  for \rhemFullTVsShortSigWord{} of the six (fixed-$b$).
\item \textbf{The gain over the HAR benchmark set.} Against the HAR-set forecast of the same model class, the full
  predictor set raises the mid-quote Sharpe ratio in \rhemFullVsHarPosWord{} of the six classes, by
  \rhrAfPoolDShMin{} to \rhrAfPoolDShMax, and the gain is significant at the 5\% level for \rhemFullVsHarSigWord{} of
  the six (Ledoit--Wolf test). For the three linear classes the comparator is the benchmark forecast itself; their
  gains, \rhemDShRidgeFullVsHar,
  \rhemDShLassoFullVsHar{} and \rhemDShEnetFullVsHar{} for the ridge, lasso and elastic net, have Ledoit--Wolf
  $p$-values of \rhemPDShRidgeFullVsHar, \rhemPDShLassoFullVsHar{} and \rhemPDShEnetFullVsHar.
\item \textbf{By calendar year} (Figure~\ref{fig:rh_econ_byyear}). Every full-set forecast has a positive $\SR_y$ in
  every year, and their median exceeds the median of the \rhemHarDistinctWord{} distinct HAR-set forecasts in
  \rhemYearsFullMedAboveHarMedWord{} of the \rhemNYearsWord{} years; the short straddle's $\SR_y$ is negative in
  \rhemYearsShortNegWord{} of them.
\item \textbf{Where the profit comes from} (Appendix~\ref{app:rh_econ_decomp}).
\item \textbf{No market beta} (Appendix~\ref{app:rh_econ_beta}). Against the excess return of the S\&P~500 the
  full-set strategies have correlations of \rhebCorrFullMin{} to \rhebCorrFullMax{} and no significant CAPM beta,
  while every CAPM $\alpha$, \rhebAlphaFullMin--\rhebAlphaFullMax\% of the premium a day, is significant. Combined with
  the market in real time, they reach Sharpe ratios of the pooled daily returns of \rhebRtMvFullMin{} to
  \rhebRtMvFullMax{} against \rhebSRmRt{} for the market alone, but the gain is significant at the 5\% level for none
  of the six mean--variance combinations or of the six inverse-volatility ones.
\item \textbf{Further results} (Appendices~\ref{app:rh_econ_sharpe}--\ref{app:rh_econ_maisch}): alternative Sharpe
  ratios, position sizing, compounding, the interval-shift test, break-even spreads, forecast loss against the Sharpe
  ratio, a longer window, and the relation to intraday option-return prediction.
\end{itemize}

